% Chapter 3: Your First Working Session
\chapter{Your First Working Session}
\label{ch:first-session}

\begin{itshape}
Your CLAUDE.md is set up and your secrets are protected.
You open Claude Code, type a prompt, and\ldots where do you start?
How do you structure a request so Claude does what you mean?
How do you know when to start fresh versus keep going?
\end{itshape}

\section{The Session Loop}
\label{sec:session-loop}

A productive session with Claude follows a repeating cycle:

\begin{enumerate}
  \item \textbf{Prompt} --- describe what you want, with enough context
    for Claude to act without guessing.
  \item \textbf{Observe} --- read what Claude did. Check the diff, run the tests,
    examine the output.
  \item \textbf{Refine} --- if the result is close but not right, give specific
    feedback. If it is wrong, start fresh (\code{/clear}).
  \item \textbf{Commit} --- when the change is verified, commit it.
    One logical change per commit.
\end{enumerate}

\margintip{This is the same loop experienced developers use
when pair programming---just faster.}

The rest of this chapter teaches three essential skills for the prompt
and refine steps: how to write clear requests, how to provide useful feedback,
and how to manage the conversation so it stays productive.

% -------------------------------------------------------------------
\section{Three Prompting Basics}
\label{sec:prompting-basics}

These three techniques handle 80\% of daily interactions.
\cref{ch:prompting} covers advanced prompting in depth;
this section gives you enough to be productive today.

\subsection{Be Specific About What, Not How}

\begin{beforeafter}
\before
\begin{minted}{text}
Fix the broken tests.
\end{minted}

\after
\begin{minted}{text}
The tests in tests/auth/ are failing with a TypeError
on line 42 of test_login.py. The error is:
"NoneType has no attribute 'email'".
Diagnose the root cause and fix it.
\end{minted}
\end{beforeafter}

The ``before'' makes Claude guess which tests, what kind of failure,
and what ``fix'' means. The ``after'' gives Claude the information
it needs to act precisely.

\subsection{Tell Claude to Think Step by Step}

\marginnote[Official]{Step-by-step prompting is consistently identified as one of the
three highest-leverage techniques.\sourceurl{https://platform.claude.com/docs/en/build-with-claude/prompt-engineering/overview}}

For complex tasks, give Claude an explicit reasoning structure:

\begin{minted}{text}
Refactor the payment module:
1. First, read the current implementation and list all public methods.
2. Then, identify methods longer than 30 lines.
3. For each long method, extract helper functions with clear names.
4. Run the existing tests after each extraction to verify nothing broke.
\end{minted}

Without this structure, Claude may attempt the entire refactoring
in one pass, increasing the risk of subtle regressions.

\subsection{Include Verification Criteria}

\convergence{``Give Claude a way to verify its own work'' is Anthropic's stated
single highest-leverage practice.\sourceurl{https://docs.anthropic.com/en/docs/claude-code/best-practices}}

Always include how to check the result:

\begin{minted}{text}
Add input validation to the create_user endpoint.
Verify: pytest tests/routes/test_users.py passes,
and manually test with: curl -X POST localhost:8000/users -d '{}'
(should return 400, not 500).
\end{minted}

When Claude knows the success criteria, it can self-correct
before you ever see the result.

% -------------------------------------------------------------------
\section{Context Hygiene: The Basics}
\label{sec:context-basics}

\marginnote[Official]{Use \code{/clear} between unrelated tasks.
Use \code{/compact} for selective summarization.\sourceurl{https://docs.anthropic.com/en/docs/claude-code/best-practices}}

The context window (introduced in \cref{sec:mental-model}) is your
workspace. Like a physical desk, it becomes less effective as it
fills with unrelated materials.

\subsection{Two Commands to Know Now}

\begin{description}
  \item[\code{/clear}] Wipe context and start fresh.
    Use between unrelated tasks: after fixing a bug, before writing a feature.
  \item[\code{/compact}] Summarize the conversation while preserving
    key decisions. Less aggressive than \code{/clear};
    use when you need to continue a long task but context is getting large.
\end{description}

\subsection{The Two-Failure Rule}

\practitioner{After two failed corrections on the same issue,
\code{/clear} and write a better initial prompt.}

Each correction attempt adds noise: the original error, your correction,
Claude's apology, and its retry. After two rounds, the context is
dominated by failure patterns rather than the actual task.
Starting fresh with a better prompt costs 30 seconds and produces
better results than a third correction round.

\begin{warnbox}[When to Clear]
If you find yourself typing ``no, that's not what I meant'' for the
third time, stop. \code{/clear}. Write the prompt you wish you had
written the first time. The cost of accumulated correction context
exceeds the cost of starting fresh.
\end{warnbox}

% -------------------------------------------------------------------
\section{A Complete Example}
\label{sec:session-example}

Here is a realistic mini-session showing the loop in practice:

\begin{minted}{text}
You: Add a rate limiter to the /api/upload endpoint.
     Token bucket algorithm, 10 requests/minute per user.
     Write tests first, then implementation.
     Verify: pytest tests/middleware/test_rate_limiter.py passes.

Claude: [writes tests, then implementation, runs tests — 2 pass, 1 fails]
  "Test test_burst_limit is failing because..."

You: The burst size should be 3, not 10. Fix the test expectation.

Claude: [fixes test, reruns — all pass]
  "All 3 tests pass. Ready to commit?"

You: Yes, commit with message "feat: Add token bucket rate limiter"
\end{minted}

Notice the structure:
\begin{itemize}
  \item The initial prompt included what, how (algorithm), and verification criteria.
  \item The refinement was specific (which test, what to change).
  \item One logical change, one commit.
\end{itemize}

% -------------------------------------------------------------------
\section{Quick Reference}

\begin{itemize}
  \item Session loop: prompt $\rightarrow$ observe $\rightarrow$ refine $\rightarrow$ commit
  \item Be specific about \emph{what} you want, not \emph{how} to do it
  \item Include verification criteria in every prompt
  \item Use \code{/clear} between unrelated tasks
  \item Two-failure rule: after two corrections, start fresh
\end{itemize}

% -------------------------------------------------------------------
\begin{trybox}
Start a session with Claude Code on your current project.
Implement one small feature or fix one bug using this loop:
\begin{enumerate}
  \item Write a prompt that includes: what you want, relevant context,
    and how to verify the result.
  \item Let Claude work. Observe the result.
  \item If refinement is needed, give specific feedback (file, line, what to change).
  \item When satisfied, commit the change.
\end{enumerate}
Track how many refinement rounds it takes.
If it exceeds two, try \code{/clear} and a better initial prompt.
\end{trybox}
